\documentclass[mode = poster]{westlakethesis}

\usepackage[width  = 45, size        = custom, 
            height = 80, orientation = portrait]{beamerposter}
\linespread{1.25}
\usepackage{newtxmath, mathtools, physics3, subcaption, siunitx, microtype}
\usephxmodu{operator}
\DeclareSIUnit\angstrom{\text{\r A}}
\DeclareSIUnit\calorie {\text{cal}}
\usepackage{chemfig}
\usepackage[version = 4]{mhchem}
\usetikzlibrary{decorations.pathmorphing, calc}
\usepackage{zref-clever, multicol}

\westlakeset{
  title     = Coherent Vibrational Wavepacket Dynamics,
  author    = Mingyu Xia (\texttt{20251202247}) \texttt|
              Department of Physics \texttt| Wenjie Dou's Group/
              Mingyu Xia \texttt| Department of Physics,
  subject   = Molecular Spectroscopy (CHEM5024) Final Project,
  bibsource = reference.bib
}

\begin{document}

\begin{frame}

\maketitle

\begin{columns}[T, onlytextwidth]
  \begin{column}{.485\linewidth}
    \begin{block}{Motivation}
      \begin{itemize}
        \item \textbf{Role:} Light-sensitive molecules drive essential
        functions, including vision and molecular switches.
        \item \textbf{Mechanism:} Ultrashort laser pulses (duration $<$
        vibrational period) excite molecules, generating coherent vibrational
        wavepackets\cite{zewail2000femtochemistry}.
        \item \textbf{Goal:} Tracking wavepacket evolution
        (\zcref[S]{fig:wavepacket}) maps excited-state potential energy
        surfaces (PES) and reveals reaction coordinates.
      \end{itemize}
    \end{block}
    \begin{exampleblock}{Theoretical Model}
      Wavepacket creation relies on the superposition of energy states via
      coherent femtosecond preparation
      \begin{equation}
        \Psi(R,t) = \sum_i c_i \psi_i(R) \exp(-\iu E_it/\hbar)
        \label{eq:wavepacket}
      \end{equation}
      \vskip-6pt
      \begin{itemize}
        \item Phase terms establish quantum coherence, allowing the system to
        evolve within a specific region.
        \item Minimal wavepacket (localized) spreading validates a classical
        picture of nuclear dynamics \cite{zewail2000femtochemistry}.
      \end{itemize}
    \end{exampleblock}
    \begin{exampleblock}{Experimental Method}
      Ultrafast dynamics are observed via Transient Absorption (TA) and
      Time-Resolved Fluorescence (TRF) \cite{liebel2015principles}.
      \begin{itemize}
        \item \textbf{Pump Pulse:} Excites molecule to Franck-Condon state;
        launches wavepacket.
        \item \textbf{Probe Pulse:} Monitors temporal evolution as PES
        wavepacket oscillations generate quantum beats.
        \item \textbf{Data Analysis:} Fourier Transform (FT) algorithms extract
        vibrational frequencies.
      \end{itemize}
    \end{exampleblock}
  \end{column}
  \hfill
  \begin{column}{.485\linewidth}
    \begin{alertblock}{Conclusions \& Outlooks}
      \begin{itemize}
        \item Quantum beat damping reveals the timescales and coupling
        strengths of conical intersection dynamics.
        \item Attosecond X-ray probing and 2D spectroscopy will dissect
        vibronic coupling, enabling coherent control of reaction pathways.
      \end{itemize}
    \end{alertblock}
    \begin{figure}[htbp]
      \centering
      \includegraphics[width = .96\linewidth]{image3.png}
      \caption{Coherent, localized wave packet (de Broglie length
      $\sim\qty{0.04}\angstrom$) for a 20 fs pulse.
      Inset: Young's Double Slit Interference.}
      \label{fig:wavepacket}
    \end{figure}
  \end{column}
\end{columns}

\vskip18pt

\begin{block}{Results and Discussion}
\begin{itemize}
  \item \textbf{Azobenzene:} TA shows two decay paths: \qty{170}{\fs}
  ($S_2 \to S_1$ via planar internal conversion) and \qty{420}{\fs}
  ($S_{3,4}$ via ring-localized dynamics), with the latter reducing
  isomerization yield.
  \item \textbf{Retinal-like Chromophores:} Coherent nuclear motion drives
  11-cis to 11-trans isomerization through a twisted conical intersection within
  \qty{200}{\femto\second}, proving initial coherent dynamics drive structural
  change \cite{schoenlein1991first}.
  \item \textbf{Organic Dyes:} Rigid dyes exhibit long-lived oscillations
  (picoseconds). TRF dynamically maps the Stokes shift before solvent
  decoherence terminates the beats.
\end{itemize}
\begin{figure}[h]
    \centering
    \begin{minipage}[b]{.48\linewidth}
      \begin{subfigure}{\linewidth}
        \centering
        \includegraphics[width = .8\linewidth]{image2.png}
        \caption{Relaxation pathways for the $S_2$, $S_3$, and $S_4$
          (near-degenerate $\pi\pi^*$) states in azobenzene.}
        \label{fig:pes_azo}
      \end{subfigure}
      \begin{subfigure}{\linewidth}
        \begin{equation}
          \setchemfig{
            bond style = {line width = 2.4pt}, double bond sep = 6pt,
            atom sep   = 1.2em, atom style = {scale = .67}}
          \begin{array}{cc}
          \chemfig[bond join, baseline = (ring1)]{*6(
            --(-)=(
              -[:30]=_[:-30]
              -[:30](-[2, 1]CH_3)=_[:-30]
              -[:30](-[2, .5,,, draw = none]11)
            =_[:-30](-[0, .5,,, draw = none]12)
              -[-2]
              (-[:-150, sqrt(3)]CH_3)
              =_[:-30]
              -[-2]
              =_[:-30]NH
              -[-2, .4]-[-2, .6,,, decorate, decoration =
                {snake, amplitude = .5mm, segment length = 2mm}]
              )-(-[3, .75])(-[1, .75])-@{ring1}-)
          }
          \xlongrightarrow{h\nu} \quad &
          \chemfig[bond join, baseline = (ring2.base)]{
            *6(
            --(-)=(
              -[:30]=_[:-30]
              -[:30](-[2, 1]CH_3)=_[:-30]
              -[:30](-[2, .5,,, draw = none]11)
            =_[:-30](-[-2,.5,,, draw = none]12)
              -[:30](-[2, 1]CH_3)=_[:-30]
              -[:30]=_[:-30]
              \chembelow{N}{H}
              -[1, .4]-[1, .6,,, decorate, decoration =
                {snake, amplitude = .5mm, segment length = 2mm}]
              )-(-[3, .75])(-[1, .75])-@{ring2}-)
          }\\
          \text{11-cis rhodopsin} & \text{All-trans photoproduct}
          \end{array}
        \end{equation}
        \caption{Reaction process of 11-cis to all-trans isomerization.}
        \label{eq:rhodopsin}
      \end{subfigure}
    \end{minipage}
    \hfill
    \begin{subfigure}[b]{.48\linewidth}
      \centering
      \begin{tikzpicture}
        \node [draw, thick, inner sep = 0pt, outer sep = 0pt] (fig2c)
          {\includegraphics[width = .54\linewidth]{image1.png}};
        \node [right] at
          ($(fig2c.south east)!\fpeval{.5/13.25}!(fig2c.north east)$) {$0$};
        \node [right] at
          ($(fig2c.south east)!\fpeval{5.5/13.25}!(fig2c.north east)$) {$50$};
        \node [right] at
          ($(fig2c.south east)!\fpeval{10.5/13.25}!(fig2c.north east)$) {$100$};
        \node [below = 1ex] at (fig2c.south) {Isomerization Coordinate};
        \node [above = 1ex, rotate around = {90:(fig2c.west)}] at (fig2c.west)
          {Energy (\unit{\kilo\calorie\per\mole})};
      \end{tikzpicture}
      \caption{Ground/excited-state PES for rhodopsin isomerization along the
        nonadiabatic path.}
      \label{fig:pes_rho}
    \end{subfigure}
    \caption{Schematic PES governing ultrafast photoisomerization.}
    \label{fig:reaction_pathways}
\end{figure}
\end{block}

\gbtsetup{shortjournal}
\begin{exampleblock}{References}
  \begin{multicols}{2}
    \small \printbibliography
  \end{multicols}
\end{exampleblock}
\end{frame}

\end{document}
